\documentclass[11pt,a4paper]{article}

%================================================
% ENCODAGE ET LANGUE
%================================================
\usepackage[utf8]{inputenc}
\usepackage[T1]{fontenc}
\usepackage[french]{babel}

%================================================
% MISE EN PAGE
%================================================
\usepackage{geometry}
\geometry{margin=1.55cm}

%================================================
% PACKAGE GTOLI
%================================================
\usepackage{../styles/gtoli}

%================================================
% INFORMATIONS DU DOCUMENT
%================================================
\title{\textbf{Cours -- 2e secondaire}\\[0.2cm]
\large Préparation au CE1D -- Mathématiques}

\author{GToli}
\date{\today}

%================================================
% DOCUMENT
%================================================
\begin{document}

\maketitle

%================================================
% OBJECTIFS
%================================================

\begin{cadreobjectif}
À la fin de ce chapitre, l’élève doit être capable de :
\begin{itemize}
    \item comprendre la notion étudiée ;
    \item utiliser correctement le vocabulaire mathématique ;
    \item appliquer une méthode structurée ;
    \item résoudre un exercice de type CE1D ;
    \item expliquer sa démarche ;
    \item vérifier la cohérence de sa réponse.
\end{itemize}
\end{cadreobjectif}

%================================================
\section{Lien avec le CE1D}
%================================================

\begin{cadrebleu}{Pourquoi ce chapitre est important ?}
Ce chapitre prépare aux attendus du premier degré et aux exercices de type CE1D.

Les exercices doivent permettre de travailler :
\begin{itemize}
    \item la compréhension de l’énoncé ;
    \item le choix de la méthode ;
    \item la mobilisation des savoirs ;
    \item la justification des démarches ;
    \item la communication claire de la réponse.
\end{itemize}
\end{cadrebleu}

%================================================
\section{Notion concernée}
%================================================

\begin{cadregris}{Ce que l’on va apprendre}
Dans ce chapitre, nous allons travailler la notion suivante :

\medskip

\textbf{Notion :} \dotfill

\medskip

Cette notion sera utilisée dans des exercices proches de ceux rencontrés au CE1D.
\end{cadregris}

%================================================
\section{Prérequis}
%================================================

\begin{cadregris}{Ce qu’il faut déjà savoir faire}
Avant de commencer, je dois être capable de :
\begin{itemize}
    \item lire correctement une consigne ;
    \item repérer les données utiles ;
    \item effectuer les calculs de base ;
    \item utiliser les notations vues en classe ;
    \item rédiger une réponse courte.
\end{itemize}
\end{cadregris}

%================================================
\section{Rappel théorique}
%================================================

\begin{cadretheorie}
Cette section contient les savoirs essentiels :

\begin{itemize}
    \item définitions ;
    \item propriétés ;
    \item formules ;
    \item vocabulaire ;
    \item exemples simples.
\end{itemize}

\medskip

\textbf{À retenir :} la théorie doit permettre de reconnaître la bonne méthode dans un exercice.
\end{cadretheorie}

%================================================
\section{Méthode CE1D}
%================================================

\begin{cadremethode}
Pour résoudre un exercice de type CE1D :

\begin{enumerate}
    \item je lis la question jusqu’au bout ;
    \item je souligne les données importantes ;
    \item je repère ce que l’on me demande ;
    \item je choisis la méthode ;
    \item j’effectue les calculs ;
    \item je vérifie si ma réponse est logique ;
    \item je rédige une phrase finale.
\end{enumerate}
\end{cadremethode}

%================================================
\section{Exemple guidé}
%================================================

\begin{cadreexemple}
\textbf{Énoncé :}

Insérer ici un exemple simple représentatif d’un exercice de 2e secondaire.

\medskip

\textbf{Analyse de l’énoncé :}
\begin{itemize}
    \item Que demande-t-on ?
    \item Quelles sont les données utiles ?
    \item Quelle méthode faut-il utiliser ?
\end{itemize}

\medskip

\textbf{Résolution guidée :}

Présenter ici les étapes de résolution.
\end{cadreexemple}

%================================================
\section{Exercice de type CE1D}
%================================================

\begin{cadreenonce}
\textbf{Exercice :}

Insérer ici un exercice inspiré du format CE1D.

\medskip

\textbf{Consigne :} résoudre l’exercice en justifiant la démarche.
\end{cadreenonce}

%================================================
\section{Aide à la résolution}
%================================================

\begin{cadreaide}
Pour t’aider :
\begin{itemize}
    \item commence par identifier ce que l’on cherche ;
    \item entoure les données utiles ;
    \item écris la formule ou la méthode ;
    \item fais attention aux unités ;
    \item vérifie que ta réponse répond bien à la question.
\end{itemize}
\end{cadreaide}

%================================================
\section{Correction détaillée}
%================================================

\begin{cadresolution}
La correction doit montrer :
\begin{enumerate}
    \item l’analyse de la consigne ;
    \item le choix de la méthode ;
    \item les calculs ;
    \item la vérification ;
    \item la phrase finale.
\end{enumerate}
\end{cadresolution}

%================================================
\section{Erreurs fréquentes}
%================================================

\begin{cadreattention}
Les erreurs fréquentes en 2e secondaire sont :
\begin{itemize}
    \item répondre trop vite sans lire toute la question ;
    \item oublier une unité ;
    \item confondre deux méthodes ;
    \item ne pas justifier ;
    \item donner uniquement un calcul sans phrase finale ;
    \item ne pas vérifier si le résultat est plausible.
\end{itemize}
\end{cadreattention}

%================================================
\section{Exercices d’entraînement}
%================================================

\begin{cadreactivite}
\begin{enumerate}
    \item Exercice d’application directe.
    \item Exercice semblable à l’exemple.
    \item Exercice avec une légère variante.
    \item Exercice de type CE1D.
    \item Exercice de synthèse.
\end{enumerate}
\end{cadreactivite}

%================================================
\section{Grille de réussite CE1D}
%================================================

\begin{cadregris}{Je vérifie ma copie}
Avant de rendre mon travail, je vérifie :
\begin{itemize}
    \item[$\square$] j’ai répondu à la question posée ;
    \item[$\square$] j’ai écrit les calculs nécessaires ;
    \item[$\square$] j’ai utilisé les bonnes unités ;
    \item[$\square$] j’ai justifié ma démarche ;
    \item[$\square$] j’ai rédigé une phrase finale ;
    \item[$\square$] mon résultat est cohérent.
\end{itemize}
\end{cadregris}

%================================================
\section{Synthèse}
%================================================

\begin{cadresynthese}
À retenir :

\begin{itemize}
    \item notion importante : \dotfill
    \item méthode principale : \dotfill
    \item formule éventuelle : \dotfill
    \item erreur à éviter : \dotfill
    \item stratégie CE1D : \dotfill
\end{itemize}
\end{cadresynthese}

%================================================
\section{Auto-évaluation}
%================================================

\begin{cadregris}{Je me situe}
\begin{itemize}
    \item[$\square$] Je comprends la notion.
    \item[$\square$] Je sais reconnaître le type d’exercice.
    \item[$\square$] Je sais choisir la méthode.
    \item[$\square$] Je sais effectuer les calculs.
    \item[$\square$] Je sais rédiger ma réponse.
    \item[$\square$] Je sais vérifier mon résultat.
\end{itemize}
\end{cadregris}

\end{document}